\begin{remark}[The assumption that underwrites all of the output]%
\label{rem:gauge}
Writing Theorem~\ref{thm:kernel} down is what exposed this. An
$\mathit{op}\in\{\mathsf{avg},\mathsf{min},\mathsf{max},\mathsf{count}\}$ rollup reduces
the points Datadog \emph{stores}. For a GAUGE those are the samples Prometheus also
holds, and \texttt{avg\_over\_time(x[$\Delta$s])} is exact. For a COUNT they are
per-interval deltas, and averaging deltas is not averaging running totals --- no
\texttt{*\_over\_time} of the cumulative series computes it, and the emitted expression is
not approximately right but unrelated. The dashboard never says which, so the emission is
correct for a gauge and the emitter records \texttt{METRIC\_TYPE\_UNKNOWN} to say so. It
reports rather than refuses
because a rule \emph{can} be chosen here; for a $\mathsf{sum}$ rollup none can, and the
same code is a refusal there instead.

The condition is not a precaution. Asked of a real account, a bare query on a
\texttt{COUNT} metric returns the \emph{mean} of four stored deltas,
so Datadog applies its enforced \texttt{avg} default rather than the \texttt{.as\_count()}
its user interface appends --- averaging deltas where \texttt{avg\_over\_time} averages
running totals (\texttt{scripts\slash trial\_checks.py e3}). The deltas are $1$, $14$,
$2$ and $3$, chosen so that each candidate predicts a different reading: $\mathsf{avg}$
gives $5$, $\mathsf{sum}$ gives $20$, a $\mathsf{count}$ rollup gives $4$ and
$\mathsf{sum}$ over the reported interval gives $1$. The account returns $5$. The first
version of this probe used $1$, $2$, $3$ and $10$, whose mean is $4$ and whose count is
also $4$, and named $\mathsf{avg}$ against a rival set that did not include
$\mathsf{count}$ --- so the reading it took could not have told the two apart.

The condition was also silent until now, and it is not rare. At an earlier stage of this
work it appeared to hold of $6{,}381$ translated occurrences --- $76.2\%$ of the
translated set as it then stood, which was smaller than today's --- because Datadog's
default time aggregator is $\mathsf{avg}$ and most queries state no method at all. The remaining quarter looked exempt: a $\mathsf{sum}$ rollup needs a
counter, the ways to ask for one are \texttt{.as\_count()}, \texttt{.as\_rate()} and
\texttt{.rollup(sum, $\Delta$)}, and all three are documented for \texttt{COUNT}
metrics --- so a query using one, we reasoned, establishes the type it does not state.

It does not, and the error was of a kind this paper has made before: a rule read in a
manual and never asked of the system. Submitting a \texttt{GAUGE} and querying it with
each modifier returns the plain rollup, unchanged and undivided --- Datadog accepts both
on a gauge and silently ignores them rather than rejecting the query
(\texttt{scripts\slash as\_rate\_order.py}). The modifier therefore records what the
dashboard's author took the metric to be. It is evidence about a person, not about a
type, and no reading of the query distinguishes the two.

The third way is weaker still, and for a different reason. \texttt{.rollup(sum, $\Delta$)}
over a \texttt{GAUGE} is not ignored: the same experiment reads it as the sum of the
stored samples, which is a perfectly meaningful query that a dashboard author might well
intend. So where \texttt{.as\_count()} on a gauge is a mistake nobody benefits from,
summing a gauge over a window is an ordinary thing to want --- and the translation this
front end chooses for it, the cumulative difference, computes something else entirely.
Of the three pieces of evidence the query can carry, the one that is hardest to dismiss
as author error is the one whose misreading is least benign. With that premise withdrawn the
caveat holds of all $8{,}475$ translated occurrences, $100\%$ of them
(Table~\ref{tab:translation-caveats}).

Recording it does not move $\kappa$ --- the translations are unchanged --- which is the
point: it changes what they are claimed to be. Its refusing counterpart still never fires,
but for a weaker reason than the one we gave: not because the method establishes the type,
but because it is the only syntax that asks for the rule, so the emitter always has a rule
to choose and never has to stop. So the same unstated fact is a diagnostic that cannot
occur and a precondition that always does.
\end{remark}
